\documentclass[11pt]{article}
\usepackage{fontspec,amsmath,amssymb,amsthm,hyperref}
\usepackage[margin=1in]{geometry}
\setlength{\emergencystretch}{3em}
\title{Bounded-test estimates from the exact corrected measure}
\author{LangevinLean mathematical exposition}
\date{}
\begin{document}\maketitle
\noindent Source module: \texttt{KineticLangevin/HamiltonianTests.lean}. Compile this proof document with XeLaTeX or LuaLaTeX. The accompanying HTML pairs the mathematical argument with the exact Lean statements and records the scope of the formalization.


The rejected move can reverse a velocity of order one. A weak approximation therefore uses bounded complex-valued tests rather than assuming the rejected jump is proportional to the step. The exact corrected measure identity implies that one uncorrected proposal has test bias at most twice the test bound times its mean rejection probability.


\section{Keep all acceptance and rejection mass}


For measurable \(0\le a\le1\), the two weighted measures \(a\pi\) and \((1-a)\pi\) sum exactly to \(\pi\). Lean combines the two densities, uses nonnegativity to rewrite the sum of \texttt{ENNReal.ofReal} values as the value of their sum, and reduces the resulting density to one. In particular both weighted measures are bounded above by the original finite measure.



\noindent Lean declarations: \texttt{acceptance\_measure\_split}.

\section{Turn the explicit invariant measure into an integral identity}


Suppose the already constructed corrected measure satisfies \(V_\#(a\pi)+R_\#((1-a)\pi)=\pi\). This identity is proved for the actual Verlet and reversal maps in the preceding module. For bounded measurable complex \(F\), integration gives \[\int F\,d\pi=\int a(z)F(Vz)\,d\pi+\int(1-a(z))F(Rz)\,d\pi.\] Lean first proves finiteness of each weighted measure from the exact mass split. The bound on \(F\) supplies integrability for both pushforward measures. It then applies the Bochner integral formula for a pushforward and the formula for a density-weighted integral; the acceptance and rejection factors are finite nonnegative real numbers.



\noindent Lean declarations: \texttt{correctedMeasure\_test\_identity}.

\section{Bound the entire rejected move through the test bound}


Subtract the invariant-test identity from the ordinary proposal expectation. The exact difference is \[\int F(Vz)\,d\pi-\int F\,d\pi=\int(1-a(z))\,[F(Vz)-F(Rz)]\,d\pi.\] If \(\|F\|\le C\), the norm of the bracket is at most \(2C\), irrespective of the size of the rejected physical jump. Since \(1-a\ge0\), the norm of the integral is at most \(2C\int(1-a)\,d\pi\). Lean explicitly checks bounded integrability of the proposal test, accepted test, rejected test, and rejection factor, rewrites the subtraction under the integral, applies the norm-of-integral inequality, and integrates the pointwise triangle bound. Thus the velocity flip is fully accounted for.



\noindent Lean declarations: \texttt{correctedMeasure\_weak\_error}.

\end{document}
